%%
%% This is file `sample-sigconf-biblatex.tex',
%% generated with the docstrip utility.
%%
%% The original source files were:
%%
%% samples.dtx  (with options: `all,proceedings,sigconf-biblatex')
%% 
%% IMPORTANT NOTICE:
%% 
%% For the copyright see the source file.
%% 
%% Any modified versions of this file must be renamed
%% with new filenames distinct from sample-sigconf-biblatex.tex.
%% 
%% For distribution of the original source see the terms
%% for copying and modification in the file samples.dtx.
%% 
%% This generated file may be distributed as long as the
%% original source files, as listed above, are part of the
%% same distribution. (The sources need not necessarily be
%% in the same archive or directory.)
%%
%%
%% Commands for TeXCount
%TC:macro \cite [option:text,text]
%TC:macro \citep [option:text,text]
%TC:macro \citet [option:text,text]
%TC:envir table 0 1
%TC:envir table* 0 1
%TC:envir tabular [ignore] word
%TC:envir displaymath 0 word
%TC:envir math 0 word
%TC:envir comment 0 0
%%
%% The first command in your LaTeX source must be the \documentclass
%% command.
%%
%% For submission and review of your manuscript please change the
%% command to \documentclass[manuscript, screen, review]{acmart}.
%%
%% When submitting camera ready or to TAPS, please change the command
%% to \documentclass[sigconf]{acmart} or whichever template is required
%% for your publication.
%%
%%
\documentclass[sigconf,natbib=false]{acmart}
%%
%% \BibTeX command to typeset BibTeX logo in the docs
\AtBeginDocument{%
  \providecommand\BibTeX{{%
    Bib\TeX}}}

%% Rights management information.  This information is sent to you
%% when you complete the rights form.  These commands have SAMPLE
%% values in them; it is your responsibility as an author to replace
%% the commands and values with those provided to you when you
%% complete the rights form.

\setcopyright{acmlicensed}
%\setcopyright{cc}
%\setcctype[4.0]{by}

\copyrightyear{2026}
\acmYear{2026}
\acmDOI{XXXXXXX.XXXXXXX}
\acmConference[SAC'26]{The 41st ACM/SIGAPP Symposium on Applied Computing}{March 23--27, 2026}{Thessaloniki, Greece}
\acmISBN{979-X-XXXX-XXXX-X/26/03}


%%
%% Submission ID.
%% Use this when submitting an article to a sponsored event. You'll
%% receive a unique submission ID from the organizers
%% of the event, and this ID should be used as the parameter to this command.
%%\acmSubmissionID{123-A56-BU3}

%%
%% For managing citations, it is recommended to use bibliography
%% files in BibTeX format.
%%
%% You can then either use BibTeX with the ACM-Reference-Format style,
%% or BibLaTeX with the acmnumeric or acmauthoryear sytles, that include
%% support for advanced citation of software artefact from the
%% biblatex-software package, also separately available on CTAN.
%%
%% Look at the sample-*-biblatex.tex files for templates showcasing
%% the biblatex styles.
%%


%%
%% The majority of ACM publications use numbered citations and
%% references, obtained by selecting the acmnumeric BibLaTeX style.
%% The acmauthoryear BibLaTeX style switches to the "author year" style.
%%
%% If you are preparing content for an event
%% sponsored by ACM SIGGRAPH, you must use the acmauthoryear style of
%% citations and references.
%%
%% Bibliography style
  \RequirePackage[
    datamodel=acmdatamodel,
    style=acmnumeric,
    backend=bibtex
    ]{biblatex}
%% Declare bibliography sources (one \addbibresource command per source)
\addbibresource{software.bib}
\addbibresource{sample-base.bib}

%%
%% end of the preamble, start of the body of the document source.
\begin{document}
%%
%% The "title" command has an optional parameter,
%% allowing the author to define a "short title" to be used in page headers.
\title{DATASUS}
% Please make sure that the short title does not exceed the width of one column
\renewcommand{\shorttitle}{Hope it is}

%%
%% The "author" command and its associated commands are used to define
%% the authors and their affiliations.
%% Of note is the shared affiliation of the first two authors, and the
%% "authornote" and "authornotemark" commands
%% used to denote shared contribution to the research.
\author{Ben Trovato}
\authornote{Both authors contributed equally to this research.}
\email{trovato@corporation.com}
\orcid{1234-5678-9012}
\author{G.K.M. Tobin}
\authornotemark[1]
\email{webmaster@marysville-ohio.com}
\affiliation{%
  \institution{Institute for Clarity in Documentation}
  \city{Dublin}
  \state{Ohio}
  \country{USA}
}

\author{Lars Th{\o}rv{\"a}ld}
\affiliation{%
  \institution{The Th{\o}rv{\"a}ld Group}
  \city{Hekla}
  \country{Iceland}}
\email{larst@affiliation.org}


\author{Julius P. Kumquat}
\affiliation{%
  \institution{The Kumquat Consortium}
  \city{New York}
  \country{USA}}
\email{jpkumquat@consortium.net}

%This command displays author info in page headers
% Please use the following convention:
% One author: J. Smith
% Two authors: J. Smith and I. Jones
% Three and more authors: J. Smith et al.
\renewcommand{\shortauthors}{B. Trovato et al.}

%%
%% The abstract is a short summary of the work to be presented in the
%% article.
\begin{abstract}
  RESUMO
\end{abstract}

%%
%% The code below is generated by the tool at http://dl.acm.org/ccs.cfm.
%% Please copy and paste the code instead of the example below.
%%
\begin{CCSXML}
<ccs2012>
 <concept>
  <concept_id>00000000.0000000.0000000</concept_id>
  <concept_desc>Do Not Use This Code, Generate the Correct Terms for Your Paper</concept_desc>
  <concept_significance>500</concept_significance>
 </concept>
 <concept>
  <concept_id>00000000.00000000.00000000</concept_id>
  <concept_desc>Do Not Use This Code, Generate the Correct Terms for Your Paper</concept_desc>
  <concept_significance>300</concept_significance>
 </concept>
 <concept>
  <concept_id>00000000.00000000.00000000</concept_id>
  <concept_desc>Do Not Use This Code, Generate the Correct Terms for Your Paper</concept_desc>
  <concept_significance>100</concept_significance>
 </concept>
 <concept>
  <concept_id>00000000.00000000.00000000</concept_id>
  <concept_desc>Do Not Use This Code, Generate the Correct Terms for Your Paper</concept_desc>
  <concept_significance>100</concept_significance>
 </concept>
</ccs2012>
\end{CCSXML}

\ccsdesc[500]{Do Not Use This Code~Generate the Correct Terms for Your Paper}
\ccsdesc[300]{Do Not Use This Code~Generate the Correct Terms for Your Paper}
\ccsdesc{Do Not Use This Code~Generate the Correct Terms for Your Paper}
\ccsdesc[100]{Do Not Use This Code~Generate the Correct Terms for Your Paper}

%%
%% Keywords. The author(s) should pick words that accurately describe
%% the work being presented. Separate the keywords with commas.
\keywords{Database, Chatbot, ?????}


%%
%% This command processes the author and affiliation and title
%% information and builds the first part of the formatted document.
\maketitle

%%%%%%%%%%%%%%%%%%%%%%%%%%%%%%%%%%%%%%%%%%%%%%%%%%%%
\section{Introduction}

% Contextualização

% Motivação 

% Objetivo e contribuições
The main goal of this work is ....

Our main contributions are outlined as follows:
\begin{itemize}
  \item BD...
  \item Agente...
  \item Dashboard...
  \item A case study demonstrating .....
\end{itemize}  


% Parágrafo com organização do texto
The remainder of this paper is structured as follows. Sections~\ref{sec:back} and~\ref{sec:related} present the background and related work. The proposed solution is described in Section~\ref{sec:solution}. ... The last section outlines our contributions, limitations, and future directions.


%%%%%%%%%%%%%%%%%%%%%%%%%%%%%%%%%%%%%%%%%%%%%%%%%%%%
\section{Background}

% Um parágrafo para contextualizar (incluir referências)
The Department of Health Information and Informatics of the Unified Health System (DATASUS)\footnote{https://datasus.saude.gov.br/}, part of the Secretariat of Digital Health and Information of the Ministry of Health (SEIDIGI/MS), was created in 1991. Its mission is to promote accessibility, integration, and security of health information, thereby strengthening the efficiency of SUS through the use of technology. Since its foundation, the department has been responsible for developing and maintaining a wide range of digital solutions that support the modernization of health management, expand connectivity in healthcare units, and improve interoperability across systems. Beyond its technological role, DATASUS also maintains and provides access to some of Brazil’s most important health databases, including the Hospital Information System (SIH/SUS), the Mortality Information System (SIM), and the National Registry of Health Establishments (CNES).

% Parágrafo para detalhar DATASUS
The DATASUS infrastructure includes data centers in Brasília and Rio de Janeiro, as well as cloud-based solutions and technological support distributed across 26 state superintendencies. Among its main initiatives are \textit{Meu SUS Digital}, \textit{SUS Profissional}, and \textit{SUS Gestor}, which connect citizens, health professionals, and managers to the system. Another key initiative is the National Health Data Network (RNDS), which enables the secure exchange of clinical information across SUS systems. At the institutional level, DATASUS formulates, implements, and monitors the National Policy on Health Information and Informatics, manages access to SUS databases, proposes semantic standards for health information, and coordinates the development of national health information systems, as defined by Decree No. 12.036 of May 28, 2024.

% Parágrafo para detalhar SIH/SUS
The Hospital Information System of the Unified Health System (SIH/SUS)\footnote{https://datasus.saude.gov.br/acesso-a-informacao/morbidade-hospitalar-do-sus-sih-sus/} was launched in 1991 by the Ministry of Health to record, process, and reimburse hospital admissions financed by SUS. It is based on the Hospital Admission Authorization (AIH), a document that consolidates administrative, clinical, and financial data for each admission in public and contracted hospitals. The SIH/SUS covers federal, state, municipal, and private hospitals accredited by SUS, making it one of the main instruments for monitoring hospital morbidity in Brazil. The database contains detailed information on patients, diagnoses (ICD-10), procedures performed, length of stay, discharge causes, and financial values. This makes it a critical source of information not only for administrative management but also for epidemiological studies.

% Último parágrafo, destacando importância e desafio
While its primary role is to enable financial reimbursement to healthcare providers, the SIH/SUS also plays an important role in supporting public policy planning, monitoring the hospital network, and guiding academic research. Due to its national scope and consistency over time, it is regarded as one of the most reliable sources of hospital morbidity data in the country. Despite this relevance, accessing SIH/SUS data still poses challenges for non-technical users. The original microdata are traditionally released in DBF (dBase) format and require extraction and transformation before being used in complex analyses. This technical barrier limits direct use by healthcare professionals and managers without advanced data processing skills, reinforcing the need for solutions that can make access easier—such as the conversational agent proposed in this work.











%%%%%%%%%%%%%%%%%%%%%%%%%%%%%%%%%%%%%%%%%%%%%%%%%%%%
\section{Related Work}
\label{sec:related}
We review prior work in three strands directly relevant to this paper: (i) conversational agents for healthcare, (ii) natural language interfaces to databases with a focus on Text-to-SQL and error repair, and (iii) healthcare-oriented Text-to-SQL applications. We purposely limit the scope here to situate our contribution within the literature; implementation details and evaluation are discussed elsewhere.

\subsection{Conversational Health Agents}
Early systematic evidence highlighted the immaturity of clinical evaluations for health chatbots and the prevalence of rule-based dialogue management \cite{laranjo2018}. More recently, agentic frameworks leverage large language models (LLMs) to orchestrate external data/tools and personalize interactions; for instance, openCHA formalizes an LLM-powered, orchestrator-centric design for health scenarios \cite{opencha2025}. In data-intensive settings, agents that generate and execute code have shown benefits for complex multi-table reasoning on electronic health records, as demonstrated by EHRAgent \cite{ehragent2024}. Meta-analytic findings further suggest that conversational agents can improve user engagement in eHealth, while underscoring the need for stronger study designs \cite{cevasco2024}.

\subsection{Natural Language Interfaces to Databases: Text-to-SQL and Repair}
Text-to-SQL has progressed from early sequence models \cite{zhong2017seq2sql} to cross-domain benchmarks and methods that explicitly model database structure. Spider established a challenging benchmark for complex, cross-domain semantic parsing \cite{yu2018spider}, followed by conversational/multi-turn variants such as SParC and CoSQL \cite{yu2019sparc,yu2019cosql}. Structural encoders (e.g., RAT-SQL) improved schema linking and contextualization \cite{ratsql2020}, while constrained decoding methods (e.g., PICARD) enforced syntactic validity at generation time \cite{picard2021}. With LLMs, surveys and empirical studies have cataloged prompting/finetuning strategies and trade-offs in accuracy and cost \cite{qin2022survey,gao2024text2sqlsurvey}. Robustness remains an active concern: Dr.Spider probes brittleness under perturbations and motivates added safeguards beyond exact-match \cite{drspider2023}. Hybrid access over structured and unstructured data has also been explored through SUQL, which extends SQL with free-text primitives \cite{suql2024}.

A complementary line addresses \emph{post-generation} correction. Clause-level edit models treat repair as operations over SQL structure and report consistent gains in exact-set-match accuracy \cite{zhang2023text}. Mutation-based strategies perform tool-agnostic repairing via candidate generation and execution feedback \cite{jiang2023repairing}. Recent agentic formulations (e.g., SQL-of-Thought; MAGIC) coordinate specialized sub-tasks—schema linking, planning, failure analysis—to guide iterative self-correction \cite{li2024sql,sun2024magic}. Comprehensive taxonomies of Text-to-SQL errors inform repair policies and evaluation protocols \cite{chen2023comprehensive}, and studies on multi-round self-debugging indicate diminishing returns beyond 1–2 iterations \cite{pourreza2024self}. Finally, in large schemas, the role of explicit schema linking is being revisited in the LLM era \cite{maamari2024death}, alongside scalable strategies for multi-database settings \cite{chen2024linkalign} and multi-LLM expert-style prompting \cite{li2024xsql}.

\subsection{Healthcare-Oriented Text-to-SQL}
Applying Text-to-SQL in healthcare introduces reliability and safety requirements that tighten tolerance for query errors. The EHRSQL shared task emphasizes reliability for EHR querying and provides targets for evaluation under realistic schema/data conditions \cite{lee2024ehrsql}. Beyond benchmark accuracy, generalization across institutions and schema variants remains difficult \cite{wang2024towards}. Graph-enhanced modeling has been proposed to better capture relational structure in medical records \cite{zhang2024graph}, while domain adaptation and knowledge integration (e.g., ICD-10, medical knowledge graphs) appear important for producing clinically meaningful queries and results \cite{kim2024clinical,roberts2024medical}.

\subsection{Conversational Systems within SUS (Brazil)}
Within the SUS ecosystem, government chatbots target service guidance (vaccination, navigation, or ombudsman intake) rather than analytical querying of national datasets \cite{ms_vax_chatbot_2023,ouvsus_chatbot_2024,meususdigital_portal}. Academic prototypes in primary care similarly emphasize addbibresource

informational support and usability—e.g., pilots with Family Health Units—without NL-to-SQL integration to DATASUS microdata \cite{jhi_aps_chatbot_2024}. To our knowledge, there are no public reports of agents that conversationally generate and validate SQL over SUS microdata (e.g., SIH/SUS via TABNET/DBF) in production settings; this positions our approach as complementary to existing informational bots and aligned with ongoing RNDS/DATASUS data-access efforts \cite{rnds_site,tabnet_manual,datasus_transfer}.


% \subsection{Research Gaps and Our Contribution}

% While significant progress has been made in text-to-SQL systems, several gaps remain in healthcare applications: (1) limited integration of domain-specific medical knowledge, (2) insufficient error correction mechanisms for safety-critical environments, and (3) lack of comprehensive frameworks combining multi-agent architectures with healthcare-specific requirements.

% Our work addresses these gaps by developing a multi-agent text-to-SQL system specifically designed for Brazilian healthcare databases, integrating medical knowledge (CID-10 classification), implementing comprehensive error correction mechanisms, and providing a production-ready framework for DATASUS data access.



\subsection{Visual Analytics in Healthcare??}


\subsection{Agents??}


%%%%%%%%%%%%%%%%%%%%%%%%%%%%%%%%%%%%%%%%%%%%%%%%%%%%
\section{Proposed Solution???}
\label{sec:solution}

% Começar falando do ambiente de desenvolvimento - ferramentas usadas. Sempre justificar as escolhas.

\subsection{Data Gathering and Pre-processing}

\subsection{Data Quality??}

\subsection{Modeling??}


\subsection{Chatbot}


\subsection{Visual Analytics}



%%%%%%%%%%%%%%%%%%%%%%%%%%%%%%%%%%%%%%%%%%%%%%%%%%%%
\section{Results/Experiments and Discussion}
\label{sec:results}



%%%%%%%%%%%%%%%%%%%%%%%%%%%%%%%%%%%%%%%%%%%%%%%%%%%%
\section{Conclusions}
\label{sec:conclusions}




%%%%%%%%%%%%%%%%%%%%%%%%%%%%%%%%%%%%%%%%%%%%%%%%%%%%
\section*{Acknowledgments}

\printbibliography[title={References}]




\end{document}
\endinput
%%
%% End of file `sample-sigconf-biblatex.tex'.
